% =============================================================================
%  Vanguard + Meridian — Deployment Runbook
%  Hand-off doc for flash through first field deploy.
% =============================================================================
\documentclass[10pt,a4paper]{article}

\usepackage[a4paper,margin=16mm,top=14mm,bottom=16mm,footskip=9mm]{geometry}

% ---- fonts ----------------------------------------------------------------
\usepackage[T1]{fontenc}
\usepackage{lmodern}
\usepackage{amssymb}
\usepackage{amsmath}
\usepackage{helvet}
\usepackage{inconsolata}
\renewcommand{\familydefault}{\sfdefault}
\usepackage{microtype}
\usepackage[none]{hyphenat}
\frenchspacing
\setlength{\parskip}{4pt}
\setlength{\parindent}{0pt}

% ---- colors ---------------------------------------------------------------
\usepackage[table]{xcolor}
\definecolor{mCyan}{HTML}{006E8F}
\definecolor{mCyanDark}{HTML}{004D64}
\definecolor{mSafe}{HTML}{2E7D32}
\definecolor{mWarn}{HTML}{E65100}
\definecolor{mFail}{HTML}{C62828}
\definecolor{mDim}{HTML}{6A7A90}
\definecolor{mInk}{HTML}{1A2332}
\definecolor{mPaper}{HTML}{FDFDFD}
\definecolor{mRow}{HTML}{F4F6FA}
\definecolor{mRule}{HTML}{CFD8E3}
\definecolor{mCode}{HTML}{0E1218}
\definecolor{mCodeInk}{HTML}{D8E0EA}
\definecolor{mAccent}{HTML}{00BFA5}

% ---- boxes ----------------------------------------------------------------
\usepackage[most]{tcolorbox}
\tcbuselibrary{breakable,skins}

% Step block: a numbered procedure card
\newtcolorbox[auto counter,number within=section]{stepbox}[1][]{
  enhanced, breakable, sharp corners=downhill,
  colback=white, colframe=mCyan, boxrule=0.6pt, arc=1pt,
  leftrule=3pt, left=8pt, right=8pt, top=8pt, bottom=7pt,
  fonttitle=\bfseries\sffamily\color{white},
  title=\MakeUppercase{Step \thetcbcounter \quad \textbullet\quad #1},
  colbacktitle=mCyan, title filled,
  attach boxed title to top left={yshift=0pt,xshift=-3pt},
  boxed title style={
    colback=mCyan, colframe=mCyan, boxrule=0pt,
    arc=0pt, left=8pt, right=8pt, top=3pt, bottom=3pt
  },
  before skip=8pt, after skip=8pt,
}

% Safety / abort callout — red, loud
\newtcolorbox{safetybox}[1][Safety]{
  enhanced, breakable,
  colback=mFail!6, colframe=mFail, boxrule=0.6pt,
  leftrule=4pt, arc=1pt, left=8pt, right=8pt, top=7pt, bottom=6pt,
  fontupper=\color{mInk},
  title={\bfseries\sffamily\color{mFail}\ding{36}\, #1},
  coltitle=mFail, colbacktitle=white, title filled=false,
  fonttitle=\bfseries\sffamily,
  before skip=6pt, after skip=6pt,
}

% Warning callout — amber
\newtcolorbox{warnbox}[1][Heads up]{
  enhanced, breakable,
  colback=mWarn!6, colframe=mWarn, boxrule=0.4pt,
  leftrule=3pt, arc=1pt, left=8pt, right=8pt, top=5pt, bottom=5pt,
  fontupper=\color{mInk},
  title={\bfseries\sffamily\color{mWarn}#1},
  coltitle=mWarn, colbacktitle=white, fonttitle=\bfseries\sffamily,
  before skip=5pt, after skip=5pt,
}

% Tip callout — cyan
\newtcolorbox{tipbox}[1][Tip]{
  enhanced, breakable,
  colback=mCyan!6, colframe=mCyan, boxrule=0.4pt,
  leftrule=3pt, arc=1pt, left=8pt, right=8pt, top=5pt, bottom=5pt,
  title={\bfseries\sffamily\color{mCyan}#1},
  coltitle=mCyan, colbacktitle=white, fonttitle=\bfseries\sffamily,
  before skip=5pt, after skip=5pt,
}

% ---- lists (enumitem) -----------------------------------------------------
\usepackage{enumitem}
\setlist[itemize]{leftmargin=*, itemsep=2pt, topsep=3pt}
\setlist[enumerate]{leftmargin=*, itemsep=2pt, topsep=3pt}
\newlist{checklist}{itemize}{1}
\setlist[checklist]{leftmargin=26pt, itemsep=4pt, topsep=0pt, label=}
\newlist{abortlist}{itemize}{1}
\setlist[abortlist]{leftmargin=22pt, itemsep=4pt, topsep=0pt, label=}

% Code block — dark monospace
\usepackage{listings}
\lstdefinelanguage{shell}{
  keywords={systemctl,sudo,tcpdump,cd,ls,pwd,mkdir,chmod,chown,rm,cp,mv,
           python,pytest,pip,cargo,git,journalctl,curl,wget},
  sensitive=true,
  morecomment=[l]{\#},
  morestring=[b]",
}
\lstset{
  language=shell,
  basicstyle=\ttfamily\small\color{mCodeInk},
  backgroundcolor=\color{mCode},
  keywordstyle=\color{mAccent}\bfseries,
  stringstyle=\color{mWarn},
  commentstyle=\color{mDim}\itshape,
  frame=none,
  xleftmargin=8pt, xrightmargin=8pt,
  columns=fullflexible, keepspaces=true,
  showstringspaces=false, tabsize=2,
  breaklines=true, breakatwhitespace=false,
  aboveskip=6pt, belowskip=6pt,
}

% ---- tables ---------------------------------------------------------------
\usepackage{booktabs}
\usepackage{array}
\usepackage{tabularx}
\newcolumntype{L}[1]{>{\raggedright\arraybackslash}p{#1}}

% ---- headings -------------------------------------------------------------
\usepackage[explicit]{titlesec}
\titleformat{\section}
  {\color{mCyan}\sffamily\bfseries\Large}
  {\llap{\color{mCyanDark}\thesection\hspace{10pt}}}
  {0pt}
  {#1}
  [\vspace{-2pt}{\color{mRule}\hrule height 0.5pt}]
\titlespacing{\section}{0pt}{16pt}{8pt}

\titleformat{\subsection}
  {\color{mCyanDark}\sffamily\bfseries\normalsize}
  {}{0pt}{#1}
\titlespacing{\subsection}{0pt}{8pt}{3pt}

% ---- hyperref -------------------------------------------------------------
\usepackage{hyperref}
\hypersetup{
  colorlinks, urlcolor=mCyan, linkcolor=mCyanDark, citecolor=mCyan,
  pdftitle={Vanguard + Meridian — Deployment Runbook},
  pdfauthor={Meridian / Vanguard},
}

% ---- symbols --------------------------------------------------------------
\usepackage{pifont}
\newcommand{\checkbox}{\raisebox{-0.4ex}{\large$\square$}\,}
\newcommand{\ding}[1]{\text{\pifont{pzd}\symbol{#1}}}
\newcommand{\badge}[2]{{\small\color{white}\colorbox{#1}{\,\strut\bfseries\sffamily #2\,}}}
\newcommand{\statusbuilt}{\badge{mSafe}{WORKING}}
\newcommand{\statustuning}{\badge{mWarn}{POLISHING}}
\newcommand{\statushw}{\badge{mDim}{NEEDS BOAT}}

% ---- fancy header & footer ------------------------------------------------
\usepackage{fancyhdr}
\pagestyle{fancy}
\fancyhf{}
\renewcommand{\headrulewidth}{0pt}
\renewcommand{\footrulewidth}{0pt}
\fancyhead[L]{\footnotesize\color{mDim}\textsf{Vanguard + Meridian $\cdot$ Deployment Runbook}}
\fancyhead[R]{\footnotesize\color{mDim}\textsf{19 April 2026}}
\fancyfoot[L]{\footnotesize\color{mDim}\textsf{Hand-off v1.1}}
\fancyfoot[C]{\footnotesize\color{mDim}\textsf{when in doubt, stop and ask}}
\fancyfoot[R]{\footnotesize\color{mDim}\textsf{\thepage}}

% ---- TOC customization ----------------------------------------------------
\usepackage{titletoc}
\titlecontents{section}[0pt]
  {\vspace{5pt}\sffamily\color{mCyanDark}\bfseries}
  {\contentslabel{20pt}}
  {\hspace{20pt}}
  {\hfill\contentspage}

\usepackage{setspace}
\usepackage{tikz}

% =============================================================================
\begin{document}
\pagestyle{empty}

% ================================ COVER ====================================
\begin{tikzpicture}[remember picture, overlay]
  \fill[mCyan] (current page.north west) rectangle
    ([yshift=-48mm]current page.north east);
  \fill[mCyanDark] ([yshift=-48mm]current page.north west) rectangle
    ([yshift=-52mm]current page.north east);
  \node[anchor=north west, xshift=18mm, yshift=-16mm, text=white]
    at (current page.north west) {
      \begin{minipage}{0.78\textwidth}
        {\Huge\bfseries\sffamily Vanguard + Meridian}\\[8pt]
        {\LARGE\sffamily Deployment Runbook}\\[4pt]
        {\normalsize\sffamily\color{white!80!mCyan}From flash to first field mission, with abort conditions}
      \end{minipage}
    };
\end{tikzpicture}

\vspace{54mm}

{\sffamily\color{mInk}\small\setstretch{1.35}
\textbf{Who this is for.} The person who will actually do the work --- flashing firmware, running the dock test, taking the boat out. It assumes familiarity with the boat itself but \emph{not} with the Meridian autopilot code. When something surprises you, stop and ask before proceeding.\par

\vspace{6pt}

\textbf{Estimated total.} Seven to ten working days from "we start flashing" to "we have field data in hand." Individual sections range from fifteen minutes (compass calibration) to several days (firmware first-boot). Don't skip §\,0.
}

\vspace{10pt}
{\color{mRule}\hrule height 0.5pt}
\vspace{10pt}

{\sffamily\large\bfseries\color{mCyanDark}Contents}
\vspace{2pt}

\begingroup
  \renewcommand{\contentsname}{}
  \tableofcontents
\endgroup

\vspace{12pt}

\begin{warnbox}[Pre-flight status --- 19 April 2026]
\setstretch{1.2}
\begin{tabular}{@{}l@{\hskip 10pt}l@{}}
Neural-net position filter: & \statusbuilt{}; three-to-five percent divergence on fine charts in harsh conditions \\
Landmark fusion (fixed 19 April): & \statusbuilt{}, with a regression test to keep it fixed \\
Parallel-filter supervisor: & \statusbuilt{} --- hot-swaps between filters on divergence \\
Cold-start, landmark-based: & \statusbuilt{} in simulation (27 m recovered) \\
Cold-start, blind search: & \statusbuilt{} --- returns a top-5 ranked candidate list; operator picks \\
Cube Plus firmware flash: & \statushw{} --- the very first step of this runbook \\
Smart motor-controller wiring: & \statushw{} --- optional; boat runs without it \\
Hours in the water on this stack: & zero --- §\,6 is the first real-hardware contact \\
\end{tabular}
\end{warnbox}

\clearpage
\pagestyle{fancy}
\setcounter{page}{1}

% ============================= §0 ===========================================
\section{Before you touch the boat}

\begin{stepbox}[Pre-flight checklist]
\begin{checklist}
\item[\checkbox] Backup of the \emph{current} autopilot firmware on the Cube Plus saved to a USB stick. \textbf{Non-negotiable.} The first boot of new firmware can surprise us, and we need a known-good restore path.
\item[\checkbox] Tablet app is paired to the long-range radio on the laptop side. Verify the link by reading heartbeat from the boat while the \emph{old} autopilot is still running. This confirms the radio layer works before we change the firmware underneath it.
\item[\checkbox] The Orca Core 2 marine gateway is flowing data on the boat's sensor network. Verify through the Orca's own web interface --- depth, satellite position, wind all reading.
\item[\checkbox] The seafloor chart tile for the deploy area is saved on the companion computer at \texttt{/opt/meridian/charts/<area>.npz}. Fetch commands are in §\,10.
\item[\checkbox] Companion computer has the \texttt{meridian-orca-bridge} service installed \textbf{but not started}.
\end{checklist}
\end{stepbox}

\begin{tipbox}[Why all this]
Every item above is a failure mode we've actually hit in lab dry-runs.
The five minutes you spend here saves four hours of troubleshooting
on the dock.
\end{tipbox}

% ============================= §1 ===========================================
\section{Flash Meridian onto the Cube Plus}
\textbf{Estimate:} 2 to 4 days. The board configuration file is drafted
but has never been compiled. Expect to find missing hardware
definitions on the first try.

\begin{stepbox}[Build the firmware image]
\begin{lstlisting}
cd <meridian repo>
./tools/build-board CubeOrangePlus --release
\end{lstlisting}
\end{stepbox}

\begin{stepbox}[Back up the current autopilot --- do not skip]
Plug the Cube into the laptop. Pull the existing firmware. Put it on
the USB stick you prepared in §\,0. If first-boot surprises us, this
restores the boat to its known-good state in ten minutes instead of an
open-ended debug session.
\end{stepbox}

\begin{stepbox}[Flash in firmware-upload mode]
Hold the BOOT button on the Cube while plugging USB, then:
\begin{lstlisting}
./tools/flash.sh CubeOrangePlus build/out/meridian.bin
\end{lstlisting}
\end{stepbox}

\begin{stepbox}[First boot --- props physically disconnected]
Watch the serial console. Expected boot sequence:
\begin{itemize}
\item Sensor/driver initialization messages
\item Sensor probe: motion unit, barometer, compass, SD card
\item Position estimator initialization
\item Radio heartbeat starts
\end{itemize}

If boot fails, diff the console output against a known-working old
autopilot log. Three most common issues:
\begin{itemize}
\item Wrong bus address for the barometer
\item Wrong pin mapping for the compass
\item SD card controller mismatch --- \textbf{fatal; must be fixed before proceeding}
\end{itemize}
\end{stepbox}

\begin{safetybox}[Do not proceed to §\,2 until you see a stable heartbeat on serial.]
\end{safetybox}

% ============================= §2 ===========================================
\section{Load the settings file}
\textbf{Estimate:} 30 minutes.

\begin{stepbox}[Load via tablet]
On the tablet: \textbf{Tools $\rightarrow$ Parameters $\rightarrow$ Load from file},
then select \texttt{vanguard\_defaults.meridian-params}.

The tablet will set every parameter one at a time. Any parameter the
firmware doesn't recognize will appear as a red "UNKNOWN PARAM" line.
We expect \textbf{zero}.
\end{stepbox}

\begin{warnbox}[If you see \emph{any} unknown parameter]
Stop. That means the settings file is out of sync with the firmware
version --- continuing would leave the boat in a half-configured
state. Send Jesse the exact parameter name.
\end{warnbox}

Expected total parameters: around 180. Confirm the \textbf{verify complete}
message before moving on.

% ============================= §3 ===========================================
\section{Bench test}
\textbf{Estimate:} 1 hour. Boat on jack stands, props free to spin with
physical clearance. Webcam pointed at the rudder so the tablet can
watch it move.

\begin{stepbox}[Arm and open Bench Test]
Arm the boat from the tablet. Confirm dialog. Open the Bench Test
overlay from the top-right menu.
\end{stepbox}

\begin{stepbox}[Rudder step test]
Use the slider, or the \textbf{STEP $\pm50$} pattern. The rudder
should follow the commanded position with less than 0.3 seconds of lag.
\end{stepbox}

\begin{stepbox}[Rudder auto-tune]
Hit \textbf{AUTOTUNE STEER}. The filter runs a step response, measures
the time constant, and suggests controller gains. \emph{Record them ---
you'll load them in a moment.}
\end{stepbox}

\begin{stepbox}[Throttle test]
Step throttle from 0 to 30\% for two seconds, then back to 0. Confirm
the prop spins and stops cleanly.
\end{stepbox}

\begin{safetybox}[Disarm before leaving the boat unattended.]
\end{safetybox}

% ============================= §4 ===========================================
\section{Compass calibration}
\textbf{Estimate:} 15 minutes. Boat on stands, or on a level dock
surface with no large iron within three metres.

\begin{stepbox}[Run the wizard]
Tablet $\rightarrow$ \textbf{Tools $\rightarrow$ Compass Cal Wizard}. Follow the
rotation sequence: nose up, nose down, roll left, roll right, full
yaw circle. The wizard writes calibration offsets to the autopilot.
Accept the summary dialog.
\end{stepbox}

% ============================= §5 ===========================================
\section{Capture the Orca's wire format}
\textbf{Estimate:} 1 day worst case --- often ten minutes.

\begin{stepbox}[Record one minute of traffic]
On the companion computer:
\begin{lstlisting}
sudo tcpdump -i eth0 -w /tmp/orca_capture.pcap
\end{lstlisting}
Wait 60 seconds with the Orca flowing its normal traffic. \texttt{Ctrl-C}
to stop.
\end{stepbox}

\begin{stepbox}[Run the autodetect]
\begin{lstlisting}
python -m meridian_orca.protocol.detect /tmp/orca_capture.pcap
\end{lstlisting}
If the output says \textbf{Yacht Devices RAW} or \textbf{Signal K},
start the bridge service and move on:
\begin{lstlisting}
systemctl --user start meridian-orca-bridge
\end{lstlisting}
\end{stepbox}

\begin{warnbox}[If autodetect says UNKNOWN]
The Orca is speaking a proprietary variant. Send the capture file to
Jesse. Expect four to eight hours to add the decoder.
\end{warnbox}

% ============================= §6 ===========================================
\section{Dock test}
\textbf{Estimate:} 2 hours. Boat in the water, tied to the slip with
at least two lines. The props may engage but cannot meaningfully
propel the boat.

\begin{stepbox}[Start the bridge and verify the sensor feed]
Start the bridge service. The tablet should show vessel traffic,
depth, wind, satellite position, and compass all flowing.
\end{stepbox}

\begin{stepbox}[Verify the depth-chart filter is live]
The status bar reads \textbf{NAV: SATELLITE} (green). The depth-chart
position marker (a dashed amber B) appears on the map within
30 seconds of the filter getting its first depth samples.
\end{stepbox}

\begin{stepbox}[Jamming-simulation test]
Press the \textbf{GPS Jam} button. Within two seconds:
\begin{itemize}
\item Badge flips to amber \textbf{DEPTH CHART (SATELLITE JAMMED)}
\item If two or more nearby vessels are broadcasting positions within
      1.5 km, the badge upgrades to cyan \textbf{DEPTH CHART + 2 LANDMARKS}
      with dashed range lines drawn to each landmark
\item The depth-chart marker takes over as the primary position indicator
\end{itemize}
Press again to restore. Badge returns to green within a few seconds.
\end{stepbox}

\begin{stepbox}[Data-freshness check]
All status-bar indicators --- satellite, depth, vessel-broadcast ---
show green dots.
\end{stepbox}

\begin{safetybox}[If anything shows red or stale, do NOT proceed to §\,7.]
\end{safetybox}

% ============================= §7 ===========================================
\section{Sheltered-water autonomous mission}
\textbf{Estimate:} 2 hours. Open water (sheltered bay or channel),
chase boat armed, manual takeover ready on the tablet.

\begin{stepbox}[Set waypoints]
Untie the boat. Long-press the map in three or four places to set a
waypoint loop.
\end{stepbox}

\begin{stepbox}[Engage AUTO]
Expected cross-track error: under 5 metres. If worse, check the
auto-tune gains from §\,3 were actually loaded.
\end{stepbox}

\begin{stepbox}[Mid-leg jamming demo]
Fire the jam button while the boat is tracking between waypoints.
The boat should continue autonomously using the depth-chart filter
for 60 seconds. Cross-track may widen to 20--30 m but the boat
should still follow the line. Restore satellite; tight tracking
resumes.
\end{stepbox}

\begin{stepbox}[Return to launch]
Trigger the return-to-launch command from the tablet. Confirm safe
return. Disarm on approach to the dock.
\end{stepbox}

% ============================= §8 ===========================================
\section{Full field deploy}
\textbf{Estimate:} 1 day on-site. Repeat §\,7 at full demo scale, with
video.

\begin{stepbox}[Capture on video]
\begin{itemize}
\item Autonomous waypoint loop (green nav)
\item Jamming segment with landmark-fusion visible (cyan badge, range lines)
\item Collision-avoidance: spawn a simulated target on collision course; watch the evasion
\item Station-keeping: drop a virtual station pin; watch the boat hold position
\end{itemize}
\end{stepbox}

\begin{stepbox}[Post-flight]
Pull telemetry logs off the companion computer through the tablet.
\end{stepbox}

% ============================= ABORT =======================================
\section{Abort conditions}
\begin{safetybox}[Stop immediately if any of these happen]
\begin{abortlist}
\item[\ding{108}] Motion-sensor vibration warning triples. Possible prop imbalance or motor-mount failure.
\item[\ding{108}] Any voltage dip below 11.0 V on a 3-cell pack, or 14.8 V on a 4-cell pack.
\item[\ding{108}] Compass heading drifts more than 30 degrees away from the satellite ground-track direction, sustained for more than five seconds.
\item[\ding{108}] Tablet loses link for more than 20 seconds. The boat \emph{should} return to launch automatically --- be ready to recover manually.
\item[\ding{108}] Water inside the hull.
\item[\ding{108}] Props spinning while the boat is marked \texttt{DISARMED}. That is a safety-interlock bug. \textbf{Kill main power.}
\end{abortlist}
\end{safetybox}

% ============================= CHARTS ======================================
\section{Chart fetching --- quick reference}

\begin{stepbox}[Public US government seafloor data (typical demo area)]
\begin{lstlisting}
python -m meridian_orca.chart_loader fetch-etopo \
  --bounds -34.00 -33.90 151.15 151.25 \
  --output /opt/meridian/charts/sydney.npz
\end{lstlisting}
\end{stepbox}

\begin{stepbox}[Global seafloor fallback (400 m detail, works anywhere)]
\begin{lstlisting}
python -m meridian_orca.chart_loader fetch-gebco \
  --bounds ... \
  --output /opt/meridian/charts/sydney_gebco.npz
\end{lstlisting}
\end{stepbox}

\begin{stepbox}[Inspect a downloaded chart]
\begin{lstlisting}
python -m meridian_orca.chart_loader info /opt/meridian/charts/sydney.npz
\end{lstlisting}
\end{stepbox}

% ============================= ROUGH EDGES ==================================
\section{Known rough edges --- current as of 19 April 2026}
\begin{itemize}
\item \textbf{Neural-net position filter} has a 3-to-5\% divergence rate on fine charts in harsh conditions. The supervisor runs the classical filter in parallel and hot-swaps on divergence, which drops end-to-end failure rate to under half a percent.
\item \textbf{Landmark fusion} had a collapse bug on 18 April (the filter would crush itself to one point and stay stuck). Fixed on 19 April with adaptive noise tolerance plus a regularization boost on fusion steps. A regression test lives in \texttt{tests/test\_ais\_fusion.py}. Validate against your own data before trusting in a new environment.
\item \textbf{Cold-start without any reference} --- no satellite, no nearby landmarks --- returns a top-5 ranked list of candidate cells. In feature-poor terrain the right answer often isn't at position \#1; an operator picks from the list or a follow-up filter converges on the one whose evidence keeps growing.
\item \textbf{Smart motor-controller upgrade} is optional. The boat runs on simple pulse-width control without it.
\item \textbf{Zero hours in the water} on this stack. §\,6 is first contact with real hardware.
\item \textbf{Remote recovery limit.} If the Cube draws power from the boat's bench/main power (not just USB), then cycling the laptop's USB controller \emph{cannot} reset the Cube. We learned this the hard way on 26 April 2026: a Phase-9 stub flashed, the chip went silent on USB, and no software-only recovery worked. Going forward, every Meridian build that touches real hardware must include USB CDC + a reboot-to-bootloader handler before deployment, so the chip is always observable and recoverable from the host. If for any reason a build is deployed without those, recovery requires a physical bench-power cycle by someone at the boat.
\end{itemize}

% ============================= SUPPORT ======================================
\section{Getting help}
For anything weird, ping Jesse. The fastest path to a fix:
\begin{enumerate}
\item \texttt{journalctl -u meridian-orca-bridge} --- last 200 lines is usually enough
\item Tablet screenshot showing the nav-source badge and status bar at the moment of the issue
\item One sentence: what you expected vs.\ what happened
\end{enumerate}

Those three things together resolve most issues inside 30 minutes.

\end{document}
